\documentclass{article}
\usepackage{graphicx} % Required for inserting images
\usepackage[russian]{babel}

\title{Мини-конспект по теме: Теорема Пифагор}
\author{Автор: Я Сам }
\date{16 июля 2025 г.}

\begin{document}


\maketitle
\tableofcontents
\newpage
\section{Введение}
Теорема Пифагора — одна из важнейших теорем евклидовой геометрии. Она находит
применение в самых разных областях:
\begin{itemize}
    \item геометрия и тригонометрия
    \item физика
    \item инженерные расчёты
    \item компьютерная графика
\end{itemize}
\section{Формулировка теоремы}
\textbf{Слова:} В прямоугольном треугольнике квадрат гипотенузы равен сумме квадратов
катетов
\begin{center}
\begin{equation}
a^2 + b^2 = c^2 \label{eq:pif}
\end{equation}
\\
Как видно из формулы \ref{eq:pif}, знание двух сторон позволяет найти третью.
\end{center}

\section{Доказательство (набросок)}
\begin{center}
Одно из доказательств основывается на площади квадрата, составленного
из четырёх одинаковых прямоугольных треугольников и малого квадрата
в центре. Раскладывая площадь двумя способами, получаем $a^2 + b^2 = c^2$.
\end{center}

\section{Пример расчёта}
\subsection{Пример 1}
\begin{center}
    a = 3, b = 4
$c=\sqrt{a^2+b^2}=\sqrt{9+16}=5$
\end{center}
\subsection{Пример 2}
\begin{enumerate}
    \item Дано: a = 5, b = 12
    \item Решение: \begin{center}
        $c=\sqrt{5^2+12^2}=\sqrt{25+144}=13$
    \end{center}
\end{enumerate}
\section{Таблица значений}
\begin{center}
    \begin{tabular}{|c|c|c|}
    \hline
    Катет \textit{a} & Катет \textit{b} & Гипотенуза \textit{c} \\
    \hline
    3 & 4 & 5 \\
    \hline
    5 & 12 & 13 \\
    \hline
    7 & 24 & 25 \\
    \hline
    \end{tabular}
\end{center}
\section{Иллюстрация}
Ниже пример изображения \\
\begin{center}
    \includegraphics[width=0.5\textwidth]{triangle.png}    
\end{center}
\section{Заключения}
Теорема Пифагора — один из краеугольных камней геометрии, помогающий решать
множество практических задач.
\section{Ссыдка и литература}
\begin{itemize}
    \item Википедия: Теорема Пифагора
    \item Классические учебники геометрии
\end{itemize}
\end{document}
